\itemhead{复习题三 · 第9题（教材p.90）}{XJ-B1-C03-RV-P090-Y1401}
\label{XJ-B1-C03-RV-P090-Y1401}
{\footnotesize 湘教版教材p.90 / PDF第97页 / 复习题三 / 第9题}
\field{原题}{已知函数 $f(x)=\dfrac{ax^2+1}{bx+c}\ (a,b,c\in\mathbb Z)$ 是奇函数，且 $f(1)=2,f(2)<3$，求 $a,b,c$ 的值。}
\field{参考答案}{$a=1,b=1,c=0$。奇函数的对称定义域与非零函数排除 $b=0$ 及非零极点的非对称情形，故 $c=0,b\ne0$；再由 $a+1=2b$、$(4a+1)/(2b)<3$ 和整数限制得唯一解。}
{\small \field{关联知识}{F08 奇偶性、对称与函数性质应用（教材帮p.113、114、115 / PDF122、123、124）；F02 定义域与原式限制（教材帮p.79、80、81 / PDF88、89、90）}
\field{方法与易错点}{检查定义域，再代入 $-x$；求另一侧解析式时同时写出 $-x$ 所在区间。 易错：只看公式不看定义域；用一对点证明奇偶；负半轴代换遗漏最外层负号。}
\field{答案依据}{答案册相关题组：PDF第24、25、26、27页（册内页23、24、25、26）；按当前题库原题逐项重算，不按旧版题号机械对套。}}
